\documentclass[11pt]{article}
\usepackage{preamble}
\usepackage{amssymb}

\newcommand{\IconCheck}{[CHECK]}
\newcommand{\IconMic}{[MIC]}
\newcommand{\IconClipboard}{[CLIPBOARD]}
\newcommand{\IconBridge}{[BRIDGE]}

\newcommand{\phasetag}[1]{%
  {\small\itshape Tag: #1\par}
}

\newcommand{\phaseitems}[1]{%
  \begin{itemize}[leftmargin=1.5em,itemsep=2pt,topsep=2pt]
  #1
  \end{itemize}%
}

\newcommand{\teacheranswerbox}[2]{%
  \begin{callout}{#1}{calloutgreen}
  #2
  \end{callout}
}

\begin{document}

\packetheading{Lesson 5-4 · Period 2 · Teacher Edition}{Radical Equations With Rational Exponents + DOK-3 Modeling}

\sectionbanner{OBJECTIVES}
{\small
\renewcommand{\arraystretch}{1.25}
\noindent\begin{tabular}{|p{1.8in}|p{4.8in}|}
\hline
\textbf{Math Objective} & Solve equations with rational exponents by converting to radical form and raising both sides to the reciprocal power; apply rational-exponent models to real-world half-life and decay problems. \\ \hline
\textbf{Language Objective} & Translate a real-world half-life scenario into an equation with a rational exponent, compute, and interpret the result in context. \\ \hline
\textbf{Essential Understanding} & Rational exponents are the bridge between algebraic power operations and real-world exponential/radical modeling — the half-life formula \(y = a\cdot(1/2)^{(t)/(h)}\) uses fractional powers directly. \\ \hline
\end{tabular}
}

\sectionbanner{TOPIC GOALS / ESSENTIAL QUESTION / MATERIALS}
{\small
\renewcommand{\arraystretch}{1.25}
\noindent\begin{tabular}{|p{1.8in}|p{4.8in}|}
\hline
\textbf{Topic Goals} & Solve rational-exponent equations; model real-world decay with fractional-power formulas. \\ \hline
\textbf{Essential Question} & How do rational-exponent equations model real-world decay, and when do extraneous solutions arise? \\ \hline
\textbf{Materials} & Student packet, pencil, calculator for Q\#41. Teacher models Ex 4; DOK-3 Practice \#41 is the topic's capstone modeling task. \\ \hline
\end{tabular}
}

\begin{callout}{\IconPin\ PRE-FLIGHT — P2 IS THE DOK-3 SPINE DAY}{calloutpink}
Today's conceptual core is Practice \#41 (Half-Life Model). Students must (a) identify \(a=50\), \(h=5\), \(t=16\), (b) compute exponent \(16/5 = 3.2\), (c) evaluate \(50\cdot(0.5)^{3.2} \approx 5.44\) mL, (d) interpret: 'about 5.44 mL of the drink remains after 16 hours.' This is the single most important item in Topic 5 — all three assessment items chain off this skill.\\
Pace: Try-It 4 (4 min) \(\rightarrow\) \#12 (4 min) \(\rightarrow\) \#33 (4 min) \(\rightarrow\) \#34 (4 min) \(\rightarrow\) \#41 (12 min). \(\sim\)35 min total Explore. (Adjust: Try-It 4 + \#12 = 8 min block, \#33 + \#34 = 8 min block.)\\
If running tight: cut \#33 or \#34. Never cut \#41.
\end{callout}

\phasetag{bridge from P1}
\frameworkphaseheader{Do Now}{1}{5}{%
\phaseitems{
  \item Hand out at door.
  \item Circulate 3 min.
  \item Cold-call 1 student for radical isolation and solution check.
}%
}{%
\phaseitems{
  \item Solve by isolating the radical and raising to the reciprocal power.
  \item Check for extraneous solutions.
}%
}{%
\phaseitems{
  \item What operation undoes a cube root?
  \item Did you verify by substituting back into the original?
}%
}{Solo teacher: circulate silently. No hints.}

\bankitem{Do Now (5-4-savvas-q21)}{%
Solve each radical equation. \(\left(\sqrt[3]{x}+8=13\right)\).
}{}

\begin{callout}{\IconTarget\ SENTENCE FRAME}{calloutyellow}
``I isolate the radical, giving \_\_\_. I raise both sides to the power \_\_\_, getting \_\_\_. I check by substituting back: the left side equals \_\_\_ and the right side equals \_\_\_. The solution is \_\_\_.''
\end{callout}

\begin{callout}{\IconCheck\ ANSWER KEY}{calloutgreen}
Answer key (from source): (x=125)
\end{callout}

\phasetag{teacher models Example 4}
\frameworkphaseheader{Launch}{2}{12}{%
\phaseitems{
  \item Project Example 4. Think aloud: rewrite the rational exponent as a radical, then raise both sides to the reciprocal power.
  \item "The exponent \(m/n\) means I raise to the \(n/m\) to undo it — the reciprocal flips the fraction."
  \item Flag the Rational-Exponent Rules card — print it in student minds.
  \item 30-sec pause: "What is the reciprocal of the exponent? Can an extraneous solution appear?" (partner check)
  \item Do NOT solve Try-Its or Practice items — release at minute 17.
}%
}{%
\phaseitems{
  \item Watch Example 4 silently.
  \item During pause: state the reciprocal exponent to partner.
  \item Copy Rules card into margin.
  \item Note the half-life formula — this comes back in \#41.
}%
}{%
\phaseitems{
  \item What must I do BEFORE raising both sides to the reciprocal power?
  \item Why do I need to check my solution even after correct algebra?
  \item In \(y = a\cdot(0.5)^{(t)/(h)}\), which value is the half-life? (\(h\))
}%
}{Project. Think aloud. Release promptly at minute 17.}

\begin{callout}{\IconBook\ RATIONAL-EXPONENT RULES (keep printed on your page)}{calloutgreen}
RATIONAL-EXPONENT RULES\\
1. Rewrite \(x^{m/n}\) as \((\sqrt[n]{x})^m\) — bridge to radicals you already know.\\
2. To solve \((\text{expr})^{m/n} = k\), raise both sides to the reciprocal power \(n/m\).\\
3. Simplify, then CHECK — even rational-exponent equations can have extraneous solutions when the denominator of the exponent is EVEN.\\
4. For real-world decay \(y = a\cdot(0.5)^{(t)/(h)}\): half-life \(h\) tells you the exponent's denominator; plug in \(t\), evaluate.
\end{callout}

\begin{callout}{\IconWarn\ EXTRANEOUS SOLUTIONS ARISE WHEN THE DENOMINATOR IS EVEN}{calloutred}
\((\text{expr})^{m/n} = k \rightarrow\) raise both sides to \(n/m\)\\
If \(n\) is EVEN, squaring can introduce solutions that don't satisfy the original equation.\\
ALWAYS substitute back into the ORIGINAL equation to verify.
\end{callout}

\bankitem{Example 4 (5-4-ex-4)}{%
{\small
Solve Equations With Rational Exponents

\textbf{A. What are the solutions to the equation \((x^2 + 5x + 5)^{5/2} = 1\)?}

\begin{align*}
(x^2 + 5x + 5)^{5/2} &= 1 \\
\left((x^2 + 5x + 5)^{5/2}\right)^{2/5} &= (1)^{2/5} \\
x^2 + 5x + 5 &= 1 \\
x^2 + 5x + 4 &= 0 \\
(x + 4)(x + 1) &= 0 \\
x + 4 = 0 &\quad \text{or} \quad x + 1 = 0 \\
x = -4 &\quad \text{or} \quad x = -1
\end{align*}

Check for extraneous solutions.
\begin{align*}
((-4)^2 + 5(-4) + 5)^{5/2} &\stackrel{?}{=} 1 &
((-1)^2 + 5(-1) + 5)^{5/2} &\stackrel{?}{=} 1 \\
(16 - 20 + 5)^{5/2} &\stackrel{?}{=} 1 &
(1 - 5 + 5)^{5/2} &\stackrel{?}{=} 1 \\
(1)^{5/2} &\stackrel{?}{=} 1 &
(1)^{5/2} &\stackrel{?}{=} 1 \\
1 &= 1 \quad \text{\textcolor{green}{\checkmark}} &
1 &= 1 \quad \text{\textcolor{green}{\checkmark}}
\end{align*}

This equation has two solutions, \(x = -4\) and \(x = -1\). There are no extraneous solutions.

\textbf{B. What is the solution to \((x + 18)^{3/2} = (x - 2)^3\)?}

\begin{align*}
(x + 18)^{3/2} &= (x - 2)^3 \\
\left((x + 18)^{3/2}\right)^{2/3} &= \left((x - 2)^3\right)^{2/3} \\
x + 18 &= (x - 2)^2 \\
x + 18 &= x^2 - 4x + 4 \\
x^2 - 5x - 14 &= 0 \\
(x + 2)(x - 7) &= 0 \\
x + 2 = 0 &\quad \text{or} \quad x - 7 = 0 \\
x = -2 &\quad \text{or} \quad x = 7
\end{align*}

Check for extraneous solutions.
\begin{align*}
(-2 + 18)^{3/2} &\stackrel{?}{=} (-2 - 2)^3 &
(7 + 18)^{3/2} &\stackrel{?}{=} (7 - 2)^3 \\
16^{3/2} &\stackrel{?}{=} (-4)^3 &
25^{3/2} &\stackrel{?}{=} 5^3 \\
64 &\neq -64 \quad \text{\textcolor{red}{X}} &
125 &= 125 \quad \text{\textcolor{green}{\checkmark}}
\end{align*}

So 7 is the only solution to the equation, and -2 is an extraneous solution.
}%
}{}

\teacheranswerbox{\IconCheck\ WORKED CONCLUSION}{%
\textbf{A.} This equation has two solutions, \(x = -4\) and \(x = -1\). There are no extraneous solutions.\\
\textbf{B.} So 7 is the only solution to the equation, and -2 is an extraneous solution.
}

\begin{callout}{\IconMic\ TEACHER SCRIPT}{calloutblue}
1. ``Watch me solve an equation with a rational exponent. I rewrite \(x^{m/n}\) as the \(n\)th root of \(x\), raised to the \(m\).''\\
2. ``Step 1: Isolate the expression with the rational exponent on one side.''\\
3. ``Step 2: Raise BOTH sides to the reciprocal power \(n/m\) to eliminate the exponent.''\\
4. ``Step 3: Simplify and solve for the variable.''\\
5. ``Step 4: CHECK by substituting back into the original — even perfect algebra can produce extraneous solutions when the denominator is even.''\\
6. ``Today's DOK-3 task applies this exact skill to a real-world half-life model. The exponent \(t/h\) IS a rational exponent.''\\
7. Release to Explore.
\end{callout}

\phasetag{TPS + DOK-3 driver}
\frameworkphaseheader{Explore}{2--3}{33}{%
\phaseitems{
  \item 4 laps across 33 min.
  \item Lap 1 (8 min): Try-It 4 + \#12 — rational-exponent equations. Press pairs to write the reciprocal power explicitly.
  \item Lap 2 (8 min): \#33 + \#34 — rational-exponent equations. Watch for pairs who skip the check step.
  \item Lap 3 (12 min): \IconStar\ \#41 Half-Life Model. Students identify \(a\), \(h\), \(t\); compute exponent; evaluate; interpret.
  \item Do NOT give \#41 answer. Pairs present computation and interpretation at Share.
}%
}{%
\phaseitems{
  \item 5 items in order: Try-It 4 \(\rightarrow\) \#12 \(\rightarrow\) \#33 \(\rightarrow\) \#34 \(\rightarrow\) \#41.
  \item For \#41: BEFORE computing, identify \(a=50\) (start amount), \(h=5\) (half-life period in hours), \(t=16\) (elapsed hours).
  \item Justify with sentence frame.
}%
}{%
\phaseitems{
  \item What is the reciprocal of the exponent? Write it before you raise both sides.
  \item \#12: does the denominator of the exponent tell you anything about extraneous solutions?
  \item \#33 / \#34: did you check your solution in the original equation?
  \item \#41: what does \(a=50\) represent? What does \(h=5\) represent?
  \item \#41: compute \(16/5\). What decimal is that? Now evaluate \((0.5)^{3.2}\).
  \item \#41: interpret your answer — what does 5.44 mL mean in context?
}%
}{Solo teacher: 4 laps. Press INTERPRETATION on \#41 — decimal answer alone is not enough.}

\textit{Work with a partner. Move in order. Practice \#41 (Half-Life Model) is the big one — plan \(\sim\)12 min for it.}

\bankitem{Try It 4 (5-4-tryit-4)}{%
Solve each equation.

\textbf{a.} \((x^2 - 3x - 6)^{3/2} - 14 = -6\)

\textbf{b.} \((x - 8)^{3/2} = (10 - x)^3\)
}{}

\begin{callout}{\IconCheck\ ANSWER / ANTICIPATED TALK}{calloutgreen}
Answer key (from source): x = 10, -7
\end{callout}

\bankitem{Practice \#12 (5-4-savvas-q12)}{%
\((3x + 2)^{2/3} = 4\)
}{}

\begin{callout}{\IconCheck\ ANSWER / ANTICIPATED TALK}{calloutgreen}
(no answer key — verify before class)
\end{callout}

\teacheranswerbox{\IconCheck\ DERIVED TEACHER ANSWER}{%
No source answer was present in the registry/docx build. Solving \((3x + 2)^{2/3} = 4\) gives \(x = 2\) or \(x = -\dfrac{10}{3}\).
}

\bankitem{Practice \#33 (5-4-savvas-q33)}{%
Solve. \(\left(0.5(x^2+5x+136)^{(2)/(3)}=50\right)\)
}{}

\begin{callout}{\IconCheck\ ANSWER / ANTICIPATED TALK}{calloutgreen}
Answer key (from source): (x=27) and (x=-32)
\end{callout}

\bankitem{Practice \#34 (5-4-savvas-q34)}{%
Solve. \(\left(2(x^2-12x-4)^{(1)/(2)}-3=15\right)\)
}{}

\begin{callout}{\IconCheck\ ANSWER / ANTICIPATED TALK}{calloutgreen}
Answer key (from source): (x=-5) and (x=17)
\end{callout}

\begin{callout}{\IconSpeak\ CHAIN 2 CALLBACK --- Half-Life (PEAK, 2 priors, cross-representational)}{calloutpeach}
\textit{When to say:} Before students start Practice \#41.

\smallskip
``Two prior versions. Venetta was polynomial: profit as a cubic. Chemist was rational: a fraction equation. Today is exponential decay: \(y = 50\cdot(0.5)^{t/5}\). The model is \textbf{given} again --- someone handed us the 5-hour half-life --- and we solve for the specific \(t\) where concentration hits a target. \textbf{Same question every time: for what input does the machine give this output?}''

\smallskip
\textit{Optional move:} after solving, ask ``what's the half-life if the function were \(y = 50\cdot(0.5)^{t/8}\)?'' --- forces students to notice the parameter is the variable that changes the machine's behavior, not the target.
\end{callout}

\bankitem{\IconStar\ DOK-3 HALF-LIFE MODEL — Practice \#41 (Topic 5 spine) (5-4-savvas-q41)}{%
Make Sense and Persevere The half-life of a certain type of soft drink is 5 h. If you drink 50 mL of this drink, the formula \(y=50(0.5)^{(t)/(5)}\) tells the amount of the drink left in your system after \(t\) hours. How much of the soft drink will be left in your system after 16 hours?
}{}

\begin{callout}{\IconStar\ DOK-3 HALF-LIFE MODEL — Practice \#41 (Topic 5 spine)}{calloutpink}
This is the single most important item in the topic. Students should (a) identify \(a=50\), \(h=5\), \(t=16\), (b) compute exponent \(16/5 = 3.2\), (c) evaluate \(50\cdot(0.5)^{3.2} \approx 5.44\) mL, (d) interpret: 'about 5.44 mL of the drink remains after 16 hours.' All three assessment items chain off this skill.
\end{callout}

\teacheranswerbox{\IconCheck\ ANSWER KEY}{%
\textbf{Answer key (from source):} about 5.44 mL
}

\phasetag{academic conversation + P3 bridge}
\frameworkphaseheader{Share / Summary}{2}{5}{%
\phaseitems{
  \item Cold-call 1 pair to present their Half-Life computation and interpretation using sentence frame.
  \item Class signals thumbs. Write the half-life formula \(y = a\cdot(0.5)^{(t)/(h)}\) on board with \(a\), \(h\), \(t\) labeled.
  \item P3 bridge: "Next period we apply these skills to assessment-critical items — bring your check-for-extraneous habit."
}%
}{%
\phaseitems{
  \item Presenting pair reads their computation and interpretation.
  \item Rest of class signals and writes takeaway in margin.
}%
}{%
\phaseitems{
  \item Summarize: what does the exponent \(t/h\) represent in the half-life formula?
  \item Does a rational exponent always produce a unique solution? (No — check for extraneous.)
}%
}{Cold-call a pair that struggled on \#41 — hold them accountable to the full interpretation step.}

\begin{sentenceframebox}
\textbf{Sentence frames:}\\
Pair share (Half-Life Model): ``This equation has exponent \_\_\_/\_\_\_, so I raise both sides to the power \_\_\_/\_\_\_. After simplifying I get \_\_\_. In the half-life formula, \(a =\) \_\_\_ (start), \(h =\) \_\_\_ (half-life period), \(t =\) \_\_\_ (elapsed time), giving \(y \approx\) \_\_\_.''\\
Whole class: one pair shares their computation and interpretation. Class signals thumbs up/sideways.
\end{sentenceframebox}

\begin{callout}{\IconBridge\ BRIDGE TO P3}{calloutpeach}
Next period we apply rational-exponent skills to assessment-critical items. Bring your check-for-extraneous habit.
\end{callout}

\phasetag{summary recap (not graded)}
\frameworkphaseheader{Exit Ticket}{1--2}{3}{%
\phaseitems{
  \item Collect at door.
  \item Scan stem 3 (50 \(\rightarrow\) 5.44 mL). Plan P3 Do Now if many students got the computation wrong or skipped the interpretation.
}%
}{%
\phaseitems{
  \item 4 sentence stems, honest.
}%
}{%
\phaseitems{
  \item Did students correctly compute \((0.5)^{3.2}\) and interpret the result in context?
}%
}{Collect. No grade.}

\begin{callout}{\IconClipboard\ EXIT TICKET — Student Form}{calloutyellow}
To solve \((\text{expr})^{m/n} = k\) I raise both sides to power \_\_\_.\\
Half-life tells me the \_\_\_ of the exponent.\\
After 16 hours, 50 mL of soft drink decays to \_\_\_ mL.\\
One biggest thing I learned today: \_\_\_\_\_\_\_\_\_\_\_\_\_\_\_\_\_\_\_\_\_\_\_\_\_\_
\end{callout}

\sectionbanner{OPTIONAL REINFORCEMENT (not collected)}
\textit{If you finish Explore early. \#19 is a conceptual extraneous-solution argument that extends the half-life skill.}

\bankitem{Optional \#1 (Savvas Practice, extraneous-solution argument)}{%
Higher Order Thinking Some, but not all, equations with rational exponents have extraneous solutions. What is the relationship between the exponents and the possibility of having extraneous solutions for equations with rational exponents? Explain your reasoning.
}{}

\teacheranswerbox{\IconCheck\ ANSWER KEY}{%
\textbf{Answer key (from source):} Only if the denominator of the rational exponent is even can there be extraneous solutions. Taking an even root of a negative number gives an extraneous solution, so if the root isn't even, there can't be an extraneous solution.
}

\clearpage

\begin{callout}{\IconClock\ PERIOD A / PERIOD F — 65-MIN CLASS NOTES}{calloutpink}
Both sections should have 65 min for P2. Use the extra 10 min on Explore \#41 (Half-Life Model) — press students to write the full interpretation sentence, not just the decimal.\\
If finished early: hand out Practice \#19 (extraneous-solution argument) as fast-finisher.
\end{callout}

\begin{callout}{\IconPuzzle\ MODIFICATIONS / SUPPORT FOR IEP STUDENTS}{calloutiep}
\textbullet\ Provide the Rational-Exponent Rules card as a sticker/cut-out.\\
\textbullet\ For \#41, provide the formula \(y = a\cdot(0.5)^{(t)/(h)}\) pre-labeled with \(a=50\), \(h=5\), \(t=16\).\\
\textbullet\ Accept verbal interpretation in lieu of written sentence.
\end{callout}

\begin{callout}{\IconGlobe\ SUPPORTS FOR ELL STUDENTS}{calloutell}
\textbullet\ Sentence frames on student page (don't skip).\\
\textbullet\ Word cards: ``reciprocal power'' = flip the exponent fraction; ``half-life'' = time for half the amount to remain; ``extraneous solution'' = an answer the algebra gives but the original equation rejects.\\
\textbullet\ ``Half-Life Model'' = a formula showing how something decreases by half every \(h\) hours.\\
\textbullet\ Pair ELL students with English-dominant partners for TPS.
\end{callout}

\end{document}
